
\documentclass[12pt]{article}
\usepackage[margin=0.75in]{geometry}
\usepackage{fontspec}
\setmainfont{Latin Modern Roman}
\usepackage{microtype}
\usepackage{multicol}
\usepackage{fancyhdr}
\usepackage{titlesec}
\usepackage{enumitem}
\usepackage{xcolor}
\usepackage[hidelinks]{hyperref}
\setlength{\columnsep}{0.28in}
\setlength{\parindent}{1.1em}
\setlength{\parskip}{0.35em}
\setlength{\headheight}{15pt}
\titleformat{\section}{\large\bfseries\scshape}{}{0pt}{}
\titleformat{\subsection}{\normalsize\bfseries}{}{0pt}{}
\pagestyle{fancy}
\fancyhf{}
\fancyfoot[C]{\thepage}
\renewcommand{\headrulewidth}{0.4pt}
\newcommand{\focusbox}[1]{\par\vspace{0.4em}\noindent\begingroup\setlength{\fboxsep}{6pt}\colorbox{gray!12}{\begin{minipage}{\dimexpr\linewidth-2\fboxsep\relax}#1\end{minipage}}\endgroup\par\vspace{0.5em}}
\newcommand{\studentline}{\vspace{0.35em}\hrule\vspace{0.65em}}

\fancyhead[L]{Whately: Argument Lab}
\fancyhead[R]{Student Handout}
\begin{document}
\begin{center}
{\LARGE\bfseries Week 1 Argument Lab}\par\vspace{0.25em}
{\large\scshape What Follows From What?}\par\vspace{0.6em}
\rule{0.82\textwidth}{0.4pt}
\end{center}

\section*{Name: \rule{2.4in}{0.4pt} \hfill Date: \rule{1.3in}{0.4pt}}

\section*{Purpose}
Whately says logic studies reasoning. Today you will practice finding the structure of an argument: the conclusion, the stated reason, the hidden assumption, and the question of whether the conclusion actually follows.

\section*{Part I: The Argument Skeleton}
For each argument, identify the parts.

\subsection*{Example}
\textbf{Argument:} "This rule is popular with most students, so it must be fair."

\begin{itemize}[leftmargin=*]
\item \textbf{Conclusion:} The rule must be fair.
\item \textbf{Stated reason:} Most students like it.
\item \textbf{Hidden assumption:} Whatever most students like is fair.
\item \textbf{Test:} The conclusion does not necessarily follow. Popularity and fairness are not the same thing.
\end{itemize}

\subsection*{1. Argument}
"Macbeth hears the witches predict his future, so he cannot be responsible for what he does afterward."

\textbf{Conclusion:}\studentline
\textbf{Stated reason:}\studentline
\textbf{Hidden assumption:}\studentline
\textbf{Does the conclusion follow? Why or why not?}\studentline\studentline

\subsection*{2. Argument}
"A person can speak confidently and still be wrong; therefore confidence is not the same thing as proof."

\textbf{Conclusion:}\studentline
\textbf{Stated reason:}\studentline
\textbf{Hidden assumption:}\studentline
\textbf{Does the conclusion follow? Why or why not?}\studentline\studentline

\subsection*{3. Argument}
"This idea is old, so it has nothing to teach modern students."

\textbf{Conclusion:}\studentline
\textbf{Stated reason:}\studentline
\textbf{Hidden assumption:}\studentline
\textbf{Does the conclusion follow? Why or why not?}\studentline\studentline

\section*{Part II: Window Notes}
Create a window note for Whately's claim that logic is useful but limited. Your window note must include:

\begin{enumerate}[leftmargin=*]
\item A drawing or visual symbol for what logic \textit{can} do.
\item A drawing or visual symbol for what logic \textit{cannot} do.
\item One key sentence from the reading.
\item A 2--3 sentence explanation of your drawing.
\end{enumerate}

\vspace{2.2in}

\section*{Part III: Exit Claim}
Complete this claim in one strong paragraph:

\focusbox{Logic is most like \rule{1.7in}{0.4pt} because \rule{2.4in}{0.4pt}.}

\vspace{1.4in}
\end{document}
